% Lösungen Skalarprodukt und Winkel (zusammenbau v0.4, Heft)
\einheitenkopf{Skalarprodukt und Winkel}

\erg{73}{a) Zahl \quad b) $\vec a \circ \vec b = 3 > 0$: spitzer Winkel \quad c) Zahl}
\erg{74}{a) (1) falsch: gleiche Richtung heißt $\vec w = k \cdot \vec v$ mit $k > 0$, die gleiche Summe erzwingt $k = 1$; (2) wahr, zum Beispiel $(600 | 0 | 0) \circ (0 | 300 | 300) = 0$ \quad b) Ausmultipliziert bleibt nur $\vec u \circ \vec u = s^2$, denn $\vec u \circ \vec v$ und $\vec u \circ \vec w$ sind null (senkrechte Kanten); h kommt nicht vor \quad c) $\vec e \circ \vec e + 3\,\vec e \circ \vec w = \lvert\vec e\rvert^2 + 0 = (s\sqrt{2})^2 = 2s^2$, weil $\vec w$ senkrecht auf der Grundfläche steht und die Diagonale $s\sqrt{2}$ lang ist}
\erg{75}{a) $\overrightarrow{AB}$ und $\overrightarrow{AC}$ \quad b) $\overrightarrow{PQ} = (2 | 1 | 2)$, $\overrightarrow{PR} = (3 | -2 | 1)$; $\mathrm{cos}\,\varphi = \frac{6}{3 \cdot \sqrt{14}}$, $\varphi \approx 57{,}7^\circ$ \quad c) $\overrightarrow{BA} \circ \overrightarrow{BC} = (-2 | -2 | -1) \circ (1 | -2 | 2) = 0$: rechter Winkel bei B}
\erg{76}{a) $\overrightarrow{KL} = (-4 | 1 | 0)$, $\overrightarrow{KS} = (-2 | -1 | 6)$; $\mathrm{cos}\,\varphi = \frac{7}{\sqrt{17} \cdot \sqrt{41}}$, $\varphi \approx 74{,}6^\circ$ \quad b) $\overrightarrow{IJ} = (-4 | 7 | -1)$, $\overrightarrow{IL} = (-6 | 0 | 6)$; $\mathrm{cos}\,\varphi = \frac{18}{\sqrt{66} \cdot \sqrt{72}}$, $\varphi \approx 74{,}9^\circ$ \quad c) $\overrightarrow{CD} \circ \overrightarrow{CS} = (-4 | 4 | 0) \circ (-4 | -4 | 3) = 0$: rechter Winkel bei C \quad d) $\overrightarrow{BC} = (-3 | 5 | 0)$; $\mathrm{cos}\,\varphi = \frac{\lvert -12 - 5 \rvert}{\sqrt{17} \cdot \sqrt{34}} = \frac{17}{17\sqrt{2}}$, also $\varphi = 45^\circ$}
\erg{77}{a) Kosinus \quad b) $\vec n = (1 | 2 | 2)$, $\lvert\vec n\rvert = 3$; $\mathrm{cos}\,\varphi = \frac{2}{3}$, $\varphi \approx 48{,}2^\circ$ \quad c) $\vec n_E = (2 | 1 | 2)$, $\vec n_F = (1 | -2 | 2)$; $\mathrm{cos}\,\varphi = \frac{4}{3 \cdot 3}$, $\varphi \approx 63{,}6^\circ$}
\erg{78}{a) $\vec n = (2 | 3 | 1)$, $\lvert\vec n\rvert = \sqrt{14}$; $\mathrm{cos}\,\varphi = \frac{1}{\sqrt{14}}$, $\varphi \approx 74{,}5^\circ$ \quad b) A, D und S haben die x-Koordinate null, F ist die Ebene $x = 0$ mit $\vec n_F = (1 | 0 | 0)$; $\mathrm{cos}\,\varphi = \frac{6}{\sqrt{85}}$, $\varphi \approx 49{,}4^\circ$ \quad c) $\mathrm{cos}\,\varphi = \frac{2}{7}$, $\varphi \approx 73{,}4^\circ$; das ist schon der verlangte Winkel, nicht das Komplement $16{,}6^\circ$ (Winkel zwischen Normalenvektor und Boden) \quad d) $\mathrm{cos}\,\varphi = \frac{4}{5}$, $\varphi \approx 36{,}9^\circ$; Winkel zwischen Dach und Wand $= 90^\circ + \varphi \approx 126{,}9^\circ$ \quad e) $\mathrm{cos}\,\varphi = \frac{8}{\sqrt{73}}$, $\varphi \approx 20{,}6^\circ$; Neigung $= \mathrm{tan}\,\varphi \cdot 100\,\% \approx 37{,}5\,\%$, die Bedingung ist erfüllt \quad f) Der Raum verschwindet, wenn die Klappe im Boden liegt; der Drehwinkel ist also der Winkel zwischen der Klappenebene und der xy-Ebene: $\mathrm{cos}\,\varphi = \frac{3}{\sqrt{13}}$, $\varphi \approx 33{,}7^\circ$ \quad g) Die Brücke liegt auf, wenn ihre Ebene die xy-Ebene ist; der Drehwinkel ist also der Winkel zwischen der Brückenebene und der xy-Ebene: $\mathrm{cos}\,\varphi = \frac{1}{\sqrt{10}}$, $\varphi \approx 71{,}6^\circ$}
\erg{79}{a) Sinus \quad b) $\lvert\vec u\rvert = 3$; $\mathrm{sin}\,\varphi = \frac{2}{3}$, $\varphi \approx 41{,}8^\circ$ \quad c) $\vec u = \overrightarrow{PS} = (-3 | -3 | 3)$, $\lvert\vec u\rvert = \sqrt{27}$; $\mathrm{sin}\,\varphi = \frac{3}{\sqrt{27}}$, $\varphi \approx 35{,}3^\circ$}
\erg{80}{a) $\lvert\vec u\rvert = \sqrt{14}$; $\mathrm{sin}\,\varphi = \frac{3}{\sqrt{14}}$, $\varphi \approx 53{,}3^\circ$}
